\section{指令滤波反步尝试}

本节尝试通过指令滤波反步来实现阶一 APF 运动
\[
    \dot{\bm p}_{i0}= -\kappa_i\bm p_{i0}+\bm\phi_i
\]
目的是避免反复微分 $\bm\phi_i$。

\subsection{理想的阶一 APF 运动}

定义
\[
    \bm\psi_i := -\kappa_i\bm p_{i0}+\bm\phi_i .
\]
若相对位置动力学可直接设为 $\dot{\bm p}_{i0}=\bm\psi_i$，
则加权 APF 能量
\[
    W(\bm p)
    :=\frac12\sum_{i=1}^{N}\kappa_i\|\bm p_{i0}\|^2+U(\bm p),
    \qquad
    \frac{\partial U}{\partial \bm p_i}=-\bm\phi_i,
\]
将满足，利用 $\sum_i\bm\phi_i=0$，
\[
\begin{aligned}
    \dot W
    &=\sum_i\kappa_i\bm p_{i0}^{T}\dot{\bm p}_{i0}
      -\sum_i\bm\phi_i^T\dot{\bm p}_i\\
    &=\sum_i\kappa_i\bm p_{i0}^{T}
      (-\kappa_i\bm p_{i0}+\bm\phi_i)
      -\sum_i\bm\phi_i^T
      (\dot{\bm p}_0-\kappa_i\bm p_{i0}+\bm\phi_i)\\
    &=-\sum_i\|\kappa_i\bm p_{i0}-\bm\phi_i\|^2\le0.
\end{aligned}
\]
因此理想的阶一运动具有直接的障碍能量论证。
由于 $U(\bm p)$ 在碰撞边界处爆破，$W$ 的有界性蕴含避碰。

\subsection{指令滤波器}

对 $m_i$ 阶相对链，定义
\[
    \bm z_{1,i}=\bm p_{i0},\quad
    \bm z_{r,i}=\bm p_{i0}^{(r-1)},\quad r=2,\dots,m_i.
\]
则
\[
    \dot{\bm z}_{r,i}=\bm z_{r+1,i},\quad r=1,\dots,m_i-1,
    \qquad
    \dot{\bm z}_{m_i,i}=\bm u_i+\bm\delta_i,
\]
其中 $\bm\delta_i=\bm d_i-\bm p_0^{(m_i)}$。

第一个指令滤波器由原始 APF 命令 $\bm\psi_i$ 驱动：
\[
    \tau_{1,i}\dot{\bm\chi}_{1,i}
    =-\bm\chi_{1,i}+\bm\psi_i,\qquad \tau_{1,i}>0.
\]
第一个实际导数的跟踪误差为
\[
    \bm e_{1,i}:=\bm z_{2,i}-\bm\chi_{1,i}.
\]
对 $r=1,\dots,m_i-2$，定义原始虚拟命令
\[
    \bar{\bm\chi}_{r+1,i}
    :=\dot{\bm\chi}_{r,i}
      -a_{r,i}\bm e_{r,i}
      -\bm e_{r-1,i},
    \qquad \bm e_{0,i}:=\bm 0,
\]
并将其通过另一个指令滤波器
\[
    \tau_{r+1,i}\dot{\bm\chi}_{r+1,i}
    =-\bm\chi_{r+1,i}+\bar{\bm\chi}_{r+1,i},
    \qquad r=1,\dots,m_i-2.
\]
对 $r=2,\dots,m_i-1$，定义
\[
    \bm e_{r,i}:=\bm z_{r+1,i}-\bm\chi_{r,i}.
\]

\subsection{标称控制律}

对 $m_i\ge2$，最后一个误差满足
\[
    \dot{\bm e}_{m_i-1,i}
    =\bm u_i+\bm\delta_i-\dot{\bm\chi}_{m_i-1,i}.
\]
一种鲁棒指令滤波反步控制候选为
\[
\begin{aligned}
    \bm u_i
    =&\ \dot{\bm\chi}_{m_i-1,i}
      -a_{m_i-1,i}\bm e_{m_i-1,i}
      -\bm e_{m_i-2,i} \\
     &-\rho_i\operatorname{sgn}(\bm e_{m_i-1,i}),
\end{aligned}
\]
其中 $a_{r,i}>0$ 且 $\rho_i>\sup_t\|\bm\delta_i(t)\|$。
若 $\bm p_0^{(m_i)}$ 已知，则项 $\bm p_0^{(m_i)}$ 可等价地加到 $\bm u_i$ 中，
并将 $\bm\delta_i$ 替换为 $\bm d_i$。

\subsection{误差估计}

令指令滤波失配为
\[
    \bm\varepsilon_{r,i}
    :=\bm\chi_{r,i}-\bar{\bm\chi}_{r,i},
    \qquad r=2,\dots,m_i-1,
\]
其中 $\bar{\bm\chi}_{1,i}:=\bm\psi_i$ 且
$\bm\varepsilon_{1,i}:=\bm\chi_{1,i}-\bm\psi_i$。
对 $r=1,\dots,m_i-2$，递归构造给出
\[
    \dot{\bm e}_{r,i}
    =-a_{r,i}\bm e_{r,i}
      -\bm e_{r-1,i}
      +\bm e_{r+1,i}
      +\bm\varepsilon_{r+1,i}.
\]
最后一个误差满足
\[
    \dot{\bm e}_{m_i-1,i}
    =-a_{m_i-1,i}\bm e_{m_i-1,i}
      -\bm e_{m_i-2,i}
      +\bm\delta_i
      -\rho_i\operatorname{sgn}(\bm e_{m_i-1,i}).
\]
因此，对
\[
    V_{e,i}:=\frac12\sum_{r=1}^{m_i-1}\|\bm e_{r,i}\|^2,
\]
交叉项按标准反步方式抵消，且
\[
\begin{aligned}
    \dot V_{e,i}
    \le&
    -\sum_{r=1}^{m_i-1}a_{r,i}\|\bm e_{r,i}\|^2
    +\sum_{r=1}^{m_i-2}\bm e_{r,i}^{T}\bm\varepsilon_{r+1,i}\\
    &-(\rho_i-\|\bm\delta_i\|)\|\bm e_{m_i-1,i}\|.
\end{aligned}
\]
因此，若指令滤波失配为零，则导数跟踪误差收敛到零。
使用有限滤波带宽时，误差关于滤波失配是输入状态稳定的，
且可通过选择小的 $\tau_{r,i}$ 使其任意小，
前提是滤波命令保持正则。

\subsection{该尝试实现了什么}

实际相对速度为
\[
    \dot{\bm p}_{i0}
    =\bm z_{2,i}
    =\bm\psi_i+\bm e_{1,i}+\bm\varepsilon_{1,i}.
\]
因此闭环实现
\[
    \dot{\bm p}_{i0}
    =-\kappa_i\bm p_{i0}+\bm\phi_i
      +\bm e_{1,i}+\bm\varepsilon_{1,i}.
\]
若 $\bm e_{1,i}\to0$ 且 $\bm\varepsilon_{1,i}\to0$，则恢复期望的
阶一 APF 运动。
若滤波失配仅较小，则系统实现一个受扰的阶一 APF 运动。

\subsection{残留难点}

该尝试避免了对 $\bm\phi_i$ 的显式高阶导数，这是指令滤波反步的主要优点。
然而，它仍不能自动给出完整的避碰证明。
沿受扰的阶一运动，
\[
    \dot{\bm p}_{i0}
    =-\kappa_i\bm p_{i0}+\bm\phi_i+\bm r_i,
    \qquad
    \bm r_i:=\bm e_{1,i}+\bm\varepsilon_{1,i},
\]
理想能量导数变为
\[
    \dot W
    =-\sum_i\|\kappa_i\bm p_{i0}-\bm\phi_i\|^2
      +\sum_i(\kappa_i\bm p_{i0}-\bm\phi_i)^T\bm r_i .
\]
最后一项可能破坏 $W$ 的单调性。
因此，要将该尝试转化为严格的围捕定理，仍需要以下额外论证之一：
\begin{itemize}
    \item 证明 $\bm r_i$ 足够小或可积，使得 $W(t)$ 对所有时间保持有界；
    \item 添加一个补偿项以恢复精确的障碍能量抵消；
    \item 添加一个强制 $\|\bm p_{ij}\|^2-d^2>0$ 的高阶 CBF 层。
\end{itemize}

因此，指令滤波反步是实现期望阶一 APF 律的有前景的实现机制，
但其本身仍不是完整的避碰证书。

\subsection{能否使 $\bm r_i$ 小或可积？}

答案是有条件的。
指令滤波构造给出标准的小增益估计，但这些估计需要 APF 命令
\[
    \bm\psi_i=-\kappa_i\bm p_{i0}+\bm\phi_i
\]
的正则性。
对于第一个指令滤波器，
\[
    \tau_{1,i}\dot{\bm\chi}_{1,i}
    =-\bm\chi_{1,i}+\bm\psi_i,\qquad
    \bm\varepsilon_{1,i}:=\bm\chi_{1,i}-\bm\psi_i,
\]
我们有
\[
    \tau_{1,i}\dot{\bm\varepsilon}_{1,i}
    =-\bm\varepsilon_{1,i}-\tau_{1,i}\dot{\bm\psi}_i.
\]
因此
\[
    \|\bm\varepsilon_{1,i}(t)\|
    \le e^{-t/\tau_{1,i}}\|\bm\varepsilon_{1,i}(0)\|
    +\tau_{1,i}\sup_{0\le s\le t}\|\dot{\bm\psi}_i(s)\|.
\]
因此，在 $\dot{\bm\psi}_i$ 有界的任何区间上，选择小的 $\tau_{1,i}$ 可使
$\bm\varepsilon_{1,i}$ 一致地小。
更强地，若 $\dot{\bm\psi}_i\in L^1[0,\infty)$ 且滤波器初始化一致，则
\[
    \|\bm\varepsilon_{1,i}\|_{L^1}
    \le \tau_{1,i}\|\dot{\bm\psi}_i\|_{L^1}.
\]
只要其输入命令足够正则，类似的估计对更高阶指令滤波器也成立。

导数跟踪误差关于滤波失配是 ISS 的。
事实上，对 $V_{e,i}$ 的估计意味着对合适的常数 $c_{1,i},c_{2,i}>0$，
\[
    \dot V_{e,i}
    \le -c_{1,i}\sum_{r=1}^{m_i-1}\|\bm e_{r,i}\|^2
       +c_{2,i}\sum_{r=2}^{m_i-1}\|\bm\varepsilon_{r,i}\|^2 .
\]
因此，小的滤波失配蕴含小的 $\bm e_{1,i}$；
平方可积的滤波失配蕴含 $\bm e_{1,i}\in L^2$。
故
\[
    \bm r_i=\bm e_{1,i}+\bm\varepsilon_{1,i}
\]
可被做小，或做成平方可积的，前提是所有滤波输入足够正则且
滤波时间常数选得足够小。

然而，这仍不是避碰证明。
$\dot{\bm\psi}_i$ 的正则性包含 $\dot{\bm\phi}_i$ 的正则性。
由于 $\bm\phi_i$ 在 $\|\bm p_{ij}\|=d$ 附近是奇异的，
$\dot{\bm\psi}_i$ 的有界性或可积性仅在安全集合的紧致子集上被保证，
即当智能体已被知远离碰撞边界时。
使用该假设来证明避碰将导致循环论证。

因此，$\bm r_i$ 的小性/可积性可作为一个局部正则结果来建立，
或在已独立证明安全裕度之后建立。
它本身不能替代前向不变性或障碍证明。